\documentclass[10pt,a4paper]{amsart}
\usepackage{amsmath,amssymb,amsthm,graphicx}
\usepackage{verbatim}
\newenvironment{reference}{\comment}{\endcomment}
%\newenvironment{reference}{}{}
\newenvironment{slogan}{\comment}{\endcomment}
\newenvironment{history}{\comment}{\endcomment}

% For commutative diagrams we use Xy-pic
\usepackage[all]{xy}

% We use 2cell for 2-commutative diagrams.
\xyoption{2cell}
% Theorem environments.
%
\theoremstyle{plain}
\newtheorem{theorem}[subsection]{Theorem}
\newtheorem{proposition}[subsection]{Proposition}
\newtheorem{lemma}[subsection]{Lemma}

\theoremstyle{definition}
\newtheorem{definition}[subsection]{Definition}
\newtheorem{example}[subsection]{Example}
\newtheorem{exercise}[subsection]{Exercise}
\newtheorem{situation}[subsection]{Situation}

\theoremstyle{remark}
\newtheorem{remark}[subsection]{Remark}
\newtheorem{remarks}[subsection]{Remarks}

\numberwithin{equation}{subsection}

% Macros
%
\def\lim{\mathop{\mathrm{lim}}\nolimits}
\def\colim{\mathop{\mathrm{colim}}\nolimits}
\def\Spec{\mathop{\mathrm{Spec}}}
\def\Hom{\mathop{\mathrm{Hom}}\nolimits}
\def\Ext{\mathop{\mathrm{Ext}}\nolimits}
\def\SheafHom{\mathop{\mathcal{H}\!\mathit{om}}\nolimits}
\def\SheafExt{\mathop{\mathcal{E}\!\mathit{xt}}\nolimits}
\def\Sch{\mathit{Sch}}
\def\Mor{\mathop{\mathrm{Mor}}\nolimits}
\def\Ob{\mathop{\mathrm{Ob}}\nolimits}
\def\Sh{\mathop{\mathit{Sh}}\nolimits}
\def\NL{\mathop{N\!L}\nolimits}
\def\CH{\mathop{\mathrm{CH}}\nolimits}
\def\proetale{{pro\text{-}\acute{e}tale}}
\def\etale{{\acute{e}tale}}
\def\QCoh{\mathit{QCoh}}
\def\Ker{\mathop{\mathrm{Ker}}}
\def\Im{\mathop{\mathrm{Im}}}
\def\Coker{\mathop{\mathrm{Coker}}}
\def\Coim{\mathop{\mathrm{Coim}}}

% Boxtimes
%
\DeclareMathSymbol{\boxtimes}{\mathbin}{AMSa}{"02}

%
% Macros for moduli stacks/spaces
%
\def\QCohstack{\mathcal{QC}\!\mathit{oh}}
\def\Cohstack{\mathcal{C}\!\mathit{oh}}
\def\Spacesstack{\mathcal{S}\!\mathit{paces}}
\def\Quotfunctor{\mathrm{Quot}}
\def\Hilbfunctor{\mathrm{Hilb}}
\def\Curvesstack{\mathcal{C}\!\mathit{urves}}
\def\Polarizedstack{\mathcal{P}\!\mathit{olarized}}
\def\Complexesstack{\mathcal{C}\!\mathit{omplexes}}
% \Pic is the operator that assigns to X its picard group, usage \Pic(X)
% \Picardstack_{X/B} denotes the Picard stack of X over B
% \Picardfunctor_{X/B} denotes the Picard functor of X over B
\def\Pic{\mathop{\mathrm{Pic}}\nolimits}
\def\Picardstack{\mathcal{P}\!\mathit{ic}}
\def\Picardfunctor{\mathrm{Pic}}
\def\Deformationcategory{\mathcal{D}\!\mathit{ef}}
\usepackage{cleveref}
\crefname{originallemma}{original lemma}{original lemmas}
\Crefname{originallemma}{Original lemma}{Original lemmas}
\makeatletter
\newlabel{algebra-lemma-field-finite-type-over-domain}{{00FY}{source}{}{sourcepointer-1}{}}
\newlabel{algebra-lemma-field-finite-type-over-domain@cref}{{[originallemma][0][]00FY}{[1][0][]source}{}{}{}}
\makeatother
\crefname{originalreference}{original reference}{original references}
\Crefname{originalreference}{Original reference}{Original references}
\makeatletter
\newlabel{lemma-rig-closed-point-in-localization}{{0GGF}{source}{}{sourcepointer-2}{}}
\newlabel{lemma-rig-closed-point-in-localization@cref}{{[originalreference][0][]0GGF}{[1][0][]source}{}{}{}}
\makeatother
\crefname{originalreference}{original reference}{original references}
\Crefname{originalreference}{Original reference}{Original references}
\makeatletter
\newlabel{lemma-rig-closed-point-relative}{{0GGB}{source}{}{sourcepointer-3}{}}
\newlabel{lemma-rig-closed-point-relative@cref}{{[originalreference][0][]0GGB}{[1][0][]source}{}{}{}}
\makeatother
\crefname{originalreference}{original reference}{original references}
\Crefname{originalreference}{Original reference}{Original references}
\makeatletter
\newlabel{lemma-rig-closed-point-relative-residue-field}{{0GGA}{source}{}{sourcepointer-4}{}}
\newlabel{lemma-rig-closed-point-relative-residue-field@cref}{{[originalreference][0][]0GGA}{[1][0][]source}{}{}{}}
\makeatother
\title{Rig-closed points survive suitable localization of the base}
\author{The Stacks Project Authors}
\date{stacks-project:restricted.tex:local-snapshot-20261012; GNU FDL 1.2 or later}
\begin{document}
\maketitle
\noindent\textit{Editorial note.} All delivered mathematical text is the author's own, copied from the pinned Stacks Project chapter file restricted.tex; no statement or proof step was rewritten, reordered or removed.\par\medskip
\noindent\textit{Editorial note.} The counted proof is the author's complete argument: the proof reduces the rig-closedness of the localized prime to a finiteness statement for the residue field, cites the finite-type residue field lemma, and applies the Nullstellensatz to choose f making the extension finite\par\medskip
\noindent\textit{Editorial note.} The label lemma-rig-closed-point-after-localization names the counted statement, and the proof environment immediately following it in the source is delivered verbatim as its proof.\par\medskip
\noindent\textit{Editorial note.} The counted statement is carried by the author's own lemma environment in the chapter source.\par\medskip
\noindent\textit{Allowed standard prerequisites (editorial).}\begin{itemize}
\item rig-closed primes and the localization B\_\{f\} of adic rings with prime q' over q (lemma-rig-closed-point-in-localization of the same chapter)
\item the equivalence characterization of rig-closed primes via finite residue field extensions (lemma-rig-closed-point-relative of the same chapter)
\item the residue field kappa of B/q is a finitely generated algebra over A/a (lemma-rig-closed-point-relative-residue-field of the same chapter)
\item the Hilbert Nullstellensatz form: a field finitely generated over a domain becomes finite after inverting a suitable element (Algebra, Lemma algebra-lemma-field-finite-type-over-domain)
\end{itemize}\medskip
\ifdefined\hypertarget\hypertarget{sourcepointer-1}{}\fi
\noindent\textit{Source reference.} Original lemma 00FY: sources/stacks-chow/tags.txt, source lines 588--588. Prerequisite pointer invoked at proof line(s) 3532; it names a standard fact of a sibling chapter that the counted argument uses, and every step of the counted proof built on it is retained verbatim in the delivered body.\par\medskip
\ifdefined\hypertarget\hypertarget{sourcepointer-2}{}\fi
\noindent\textit{Source reference.} Original reference 0GGF: sources/stacks-restricted/pinned-original-source.tex, source lines 3460--3460. Prerequisite pointer invoked at proof line(s) 3521; it names a standard fact of the same chapter that the counted argument uses, and every step of the counted proof built on it is retained verbatim in the delivered body.\par\medskip
\ifdefined\hypertarget\hypertarget{sourcepointer-3}{}\fi
\noindent\textit{Source reference.} Original reference 0GGB: sources/stacks-restricted/pinned-original-source.tex, source lines 3333--3333. Prerequisite pointer invoked at proof line(s) 3517; it names a standard fact of the same chapter that the counted argument uses, and every step of the counted proof built on it is retained verbatim in the delivered body.\par\medskip
\ifdefined\hypertarget\hypertarget{sourcepointer-4}{}\fi
\noindent\textit{Source reference.} Original reference 0GGA: sources/stacks-restricted/pinned-original-source.tex, source lines 3314--3314. Prerequisite pointer invoked at proof line(s) 3529; it names a standard fact of the same chapter that the counted argument uses, and every step of the counted proof built on it is retained verbatim in the delivered body.\par\medskip
\begin{lemma}
\label{lemma-rig-closed-point-after-localization}
Let $\varphi : A \to B$ be an arrow of $\textit{WAdm}^{Noeth}$
which is adic and topologically of finite type.
Let $\mathfrak q \subset B$ be rig-closed.
There exists an $f \in A$ which maps to a unit in
$B/\mathfrak q$ such that we obtain a diagram
$$
\vcenter{
\xymatrix{
B \ar[r] &
B_{\{f\}} \\
A \ar[r] \ar[u]_\varphi &
A_{\{f\}} \ar[u]_{\varphi_{\{f\}}}
}
}
\quad\text{with primes}\quad
\vcenter{
\xymatrix{
\mathfrak q \ar@{-}[r] \ar@{-}[d] &
\mathfrak q' \ar@{-}[d] \ar@{=}[r] &
\mathfrak q B_{\{f\}} \\
\mathfrak p \ar@{-}[r] &
\mathfrak p'
}
}
$$
such that $\mathfrak p'$ is rig-closed, i.e.,
the map $A_{\{f\}} \to B_{\{f\}}$ and the prime ideals
$\mathfrak q'$ and $\mathfrak p'$ satisfy
the equivalent conditions of Lemma \ref{lemma-rig-closed-point-relative}.
\end{lemma}

\begin{proof}
Please see Lemma \ref{lemma-rig-closed-point-in-localization}
for the description of $\mathfrak q'$. The only assertion the lemma makes
is that for a suitable choice of $f$ the prime ideal $\mathfrak p'$
has the property $\dim((A_f)^\wedge/\mathfrak p') = 1$.
By Lemma \ref{lemma-rig-closed-point-relative} this in turn
just means that the residue field $\kappa$ of
$B/\mathfrak q = (B_f)^\wedge/\mathfrak q'$ is finite over
$(A_f)^\wedge/\mathfrak a' = (A/\mathfrak a)_f$.
By Lemma \ref{lemma-rig-closed-point-relative-residue-field}
we know that $A/\mathfrak a \to \kappa$ is a finite type
algebra homomorphism. By the Hilbert Nullstellensatz in
the form of Algebra, Lemma \ref{algebra-lemma-field-finite-type-over-domain}
we can find an $f \in A$ which maps to a unit in $\kappa$
such that $\kappa$ is finite over $A_f$. This finishes the proof.
\end{proof}

\end{document}
